% ============================================================================
% SECTION 1 — INTRODUCTION
% Paper budget: ~1.5 pages. The order is fixed: context → gap → what we do →
% research questions → contributions → roadmap. Do not add subsections beyond
% the ones here; an introduction that runs to four subsections reads like a
% thesis chapter, which is exactly the note this document is fixing.
% ============================================================================
\section{Introduction}
\label{sec:introduction}

Agentic systems built on large language models (LLMs) now write and maintain
production software every day, and practitioners are still learning where they
can be trusted. Building software moves through planning, implementation,
debugging, review, and maintenance. Each phase has its own goals and its own
standard of correctness, and each fails in its own ways. A general-purpose
model that does well in one phase offers no guarantee in another, so agent
behaviour has to be studied phase by phase.

This paper is about debugging, which is expensive. Before writing any part of
a fix, a developer reads or writes a bug report, gathers diagnostic
information, often builds a test case that exposes the failure, and looks for
the root cause. The work is repetitive and takes a large share of the
development cycle, so it is an obvious target for automation, and a hard one.

LLM-based agents have made progress on code generation, test generation, and
program repair, but fixing bugs in real projects remains difficult. On
SWE-bench, a benchmark of 2,294 GitHub issues introduced by Jimenez
et al.~\cite{Jimenez2023SWEbench}, the SWE-agent system of Yang
et al.~\cite{Yang2024SWEagent} resolves about 12.5\% of issues. That was the
state of the art when it was published, and it shows how far automated repair
still is from being reliable on real code.

A \emph{bug reproduction test} (BRT) is a test that fails on the buggy
revision and passes once the bug is fixed, and BRTs help repair. In a study on
Google's internal codebase, Cheng et al.~\cite{Cheng2025Agentic} show that
giving an automated program repair system high-quality BRTs increases the
share of bugs that reach a plausible fix.

That finding has not reached agent \emph{architecture}. Current repair agents
are reactive: they write a patch first and validate it afterwards, so a test, if
one appears at all, can only confirm a patch that already exists and never
steers the search for one. The evidence for BRTs also comes from proprietary
internal issues; it has not been tested on the public benchmarks used to
compare repair agents. This paper changes the order. Before any fix is
attempted, an agent must write a reproduction test from the bug report and the
repository, and must \emph{observe} that test fail on the buggy revision. The
ordering follows test-driven development as people practise it, and it lets
the agent use its own reading of the code to sharpen the goal it then has to
meet.

Our hypothesis is that a reproduction test, as an executable specification of
the failure, narrows the set of acceptable patches. Requiring it first should
give the repair stage earlier and sharper feedback than checking a patch after
the fact, which should mean fewer unfocused edits and more precise repairs.

\subsection{Problem Statement}
\label{sec:intro-problem}

Repair agents follow no standard procedure for a defect. What an agent does
next depends on the user's prompt, and little constrains the failure it is
supposed to reproduce. A reactive agent has no executable specification of the
bug. It works from its own reading of the issue text, which leaves room to
misread the issue, to make unfocused edits, to break unrelated code, and to
decide for itself when it is done. Test-driven development addresses these
failure modes for human developers, because every change is checked against a
test that states what ``fixed'' means. LLM test generation is well studied,
but as a system of its own, outside the repair agent. The open question is
therefore architectural: \emph{does making reproduction mandatory and prior, inside a
single agentic pipeline, resolve more real-world issues than validating
reactively?}

\subsection{Research Questions}
\label{sec:intro-rqs}

How many real issues an agentic system resolves depends on the quality of its
tests, its repair capability, and its orchestration~\cite{Meng2024Empirical}.
The literature does not say how to combine these components or in which order
to run them. Our first question asks whether the ordering matters, the second
which pipeline characteristics help under a fixed budget, and the third which
models can run the pipeline at all.

\textbf{RQ1 (Ordering).} Does a test-first agentic pipeline, in which repair is
withheld until a generated test has been observed to fail on the buggy
revision, resolve more real-world issues than a matched reactive pipeline? We
answer it with a paired comparison on SWE-bench Lite. Both arms process every
sampled instance with the same model, tool set, iteration and token budgets,
container, and scoring harness, and they differ only in the reproduction phase
and the tool that runs its test. \emph{Reactive} is the usual design, as in
SWE-agent~\cite{Yang2024SWEagent}: the agent may run tests, but nothing
requires it to produce a failing one before it starts editing. The primary
outcome is resolved-rate, where an instance counts as resolved when the
benchmark's held-out failing tests pass and its previously passing tests still
pass.

\textbf{RQ2 (Orchestration).} Which pipeline and orchestration characteristics
improve fix quality and reliability under a fixed budget? An agentic pipeline
can be designed in many ways: how many agents take part, which tools each may
call, how they hand work to one another, and what evidence a test or a fix
needs before it is accepted. We ablate the test-first pipeline one component
at a time, with model, sample, container, and harness held fixed. The
components are the criterion for accepting a generated test without the
reference patch, the per-instance attempt budget, self-review of the test and
of the fix, adversarial verification after the fix, and best-of-$N$ test
selection. For each variant we report the change in resolved-rate and the
measured cost per instance. Two ablations produced negative results, and we
report them as findings.

\textbf{RQ3 (Open-weight models).} Can open-weight models sustain the
test-first pipeline, and at which stage does the trajectory break down?
Open-weight models are attractive for repair work because they are cheap to
run and can be deployed on one's own hardware. Being able to write code does
not mean a model can act as an agent. We run the same pipeline, prompts,
tools, and instances as the frontier arm, and record where each trajectory
ends: whether the agent reaches a terminal submission within budget, whether
the test it submits reproduces the failure on the buggy revision, and whether
the fix it then produces resolves the issue. Reporting where the pipeline stops,
in addition to resolved-rate, separates a limit of the scaffolding from a limit
of the model, and that distinction tells us whether better prompts could close
the gap.

\subsection{Contributions}
\label{sec:intro-contributions}

% A paper states its contributions explicitly; a thesis usually does not, so
% this list is new. It is drafted from material already in PLAN.md and is
% yours to rewrite — keep it to four items, and make each one a thing we BUILT
% or MEASURED, never a thing we "discuss". Reviewers read this list as the
% paper's claim.
This paper makes the following contributions.

\begin{itemize}
  \item \textbf{Test-first ordering as an architectural constraint.} We present
  TDBF, a two-agent pipeline in which the repair agent is not reachable until a
  generated test has been observed to fail on the buggy revision, and we
  contrast it with prompt-level test-first instructions inside a single agent.
  \item \textbf{A gold-free acceptance gate.} We define and implement a test
  acceptance criterion that never consults the reference patch, the fixed tree,
  or the held-out tests, so the decision is one a deployed system could make;
  reference data is used only offline, as a measurement instrument.
  \item \textbf{A paired evaluation on public data.} We compare test-first
  against a matched reactive baseline on SWE-bench Lite under identical model,
  tools, budgets, and scoring, reporting resolved-rate with paired significance
  testing and per-instance cost.
  \item \textbf{An orchestration ablation, including negatives.} We measure the
  acceptance criterion, attempt budget, self-review, post-fix adversarial
  verification, and best-of-$N$ selection, and report the configurations that
  did not help alongside the ones that did.
\end{itemize}

% Roadmap: one sentence, present tense, no "we will". Update if sections move.
Section~\ref{sec:related} positions the work against test generation, agentic
repair, and the one prior study of test ordering.
Section~\ref{sec:approach} documents the pipeline as built.
Section~\ref{sec:setup} gives the experimental design,
Section~\ref{sec:results} the results, and
Section~\ref{sec:discussion} their interpretation and the threats to validity.
Section~\ref{sec:conclusion} concludes.
